%! @@TITLE@@ — Unit @@UNITINT@@, Lesson @@LESSONID@@ — Guided Notes & Practice
% GRADUAL RELEASE component (I do -> we do), 34 minutes.
%   I DO   : the objective box, the vocabulary box, and EXACTLY TWO long notes sections. The
%            teacher presents; students fill the blanks as each idea is named. Each section
%            carries two moves — the second introduced by a bold run-in heading, not a new box.
%   WE DO  : the practicebox ("Guided Practice") — worked together, students hold the pen.
%            THE NOTES END HERE.
% Vocabulary is GIVEN here, not discovered.
% THERE IS NO INDEPENDENT PRACTICE SET, no "Putting It Together", and no reflectionbox. The
% HOMEWORK is the individual practice: students start it in class in the last 8 minutes, and
% that supervised start is the teacher's second formative read. There is no exit ticket.
%
% PAGE LOCKSTEP with notes_key/: every \blank{} here becomes a short \ans{} there, and every
% \writelines{n} becomes exactly n \ansline{} calls. SIZE EACH \blank{} CLOSE TO ITS ANSWER —
% a 4.4cm blank replaced by a 2cm answer rewraps the paragraph and costs a page. Every
% \begin{work} block is authored BYTE-IDENTICALLY in both files.
% VOCABPAR: \par brackets each \termblank/\termblanklong inside the vocabbox.
\documentclass[12pt]{article}
\usepackage{@@PREFIX@@-article}
\usepackage{@@PREFIX@@-boxes}

\begin{document}

\pageheader{Unit @@UNITINT@@, Lesson @@LESSONID@@}{Guided Notes \& Practice}
% NAMESTRIP: no \namedateperiod here — the name row lives on the cover only.

\vspace{0.12in}

\boxguard[14]
\begin{objectivebox}
\small
By the end of these notes I will be able to:
\begin{enumerate}[leftmargin=1.6em, nosep, after=\vspace{2pt}]
  \item TODO: target 1, in student language, naming the formal term;
  \item TODO: target 2;
  \item TODO: target 3.
\end{enumerate}
\end{objectivebox}

\vspace{0.10in}

% ── VOCABULARY (I do — students fill each term AS IT IS NAMED, never in advance) ──
\boxguard[20]
\begin{vocabbox}
\small
Fill in each term \textbf{as we name it} in the notes below.
\par
\termblank{TODO Term 1}
\par
\termblank{TODO Term 2}
\par
\termblank{TODO Term 3}
\par
\end{vocabbox}

\vspace{0.10in}

% ── I DO: EXACTLY TWO long notes sections (~20 minutes total) ────────────────
% The crux — the item that surfaces the lesson's target misconception — lives in the
% second half of section 2.
% Each section: a short piece of exposition, a PRE-DRAWN display or table, and \blank{} fills
% that students complete as the idea is named. Never "sketch a graph" — every dotplot,
% histogram, boxplot, scatterplot, and computer output is GIVEN.
\boxguard[12]
\begin{notesbox}{1. TODO Section Title --- and TODO Second Move}
\small
TODO: exposition, 2--3 sentences, in a real context.

\vspace{6pt}
\textbf{TODO Formal term} is \blank{2.4cm} --- TODO definition with fills sized close to
their answers.

% The second half of section 1 — a bold run-in heading, NOT a new box.
\vspace{10pt}
\textbf{\textcolor{navy}{TODO Second Move.}}
TODO: the test, rule, or contrast that follows directly from the definition above.
\end{notesbox}

\vspace{0.10in}

\boxguard[12]
\begin{notesbox}{2. TODO Section Title --- and TODO Second Move}
\small
TODO: a pre-drawn table or display, then fills that read it.

\vspace{5pt}
\renewcommand{\arraystretch}{1.25}
\begin{tabularx}{\linewidth}{@{} l Y Y @{}}
\rowcolor{sky}
\textbf{TODO} & \textbf{TODO} & \textbf{TODO} \\
\hline
TODO & TODO & TODO \\
\hline
\end{tabularx}

\vspace{6pt}
\begin{center}
\fcolorbox{redacc}{redbg}{\parbox{0.94\linewidth}{\small
\textbf{TODO: the target misconception, called out by name.} TODO: the probe that breaks it.}}
\end{center}

% The second half of section 2 — a bold run-in heading, NOT a new box. This is the crux.
% The \begin{work} block takes NO argument and its body is math (an amsmath aligned):
% one statement per line, & immediately before the relation. Author it BYTE-IDENTICALLY in
% the key — the blank reserves its exact height and prints nothing.
\vspace{10pt}
\textbf{\textcolor{navy}{TODO Second Move.}}
TODO: the computation, modelled by the teacher.
\begin{work}
  \text{TODO statistic} &= \text{TODO formula} \\
                        &= \text{TODO substitution} \\
                        &\approx \text{TODO value}
\end{work}
Interpreted in context, this says \blank{6.0cm}.
\end{notesbox}

\vspace{0.10in}

% ── WE DO: guided practice (~6 min) — students hold the pen, teacher holds the questions ──
% practicebox takes NO argument; its title is fixed as "Guided Practice".
\boxguard[20]
\begin{practicebox}
\small
\textbf{TODO Title --- we work this one together.} TODO: one problem in a fresh context.

\begin{enumerate}[label=\textbf{\alph*.}, leftmargin=1.8em, itemsep=7pt]
  \item TODO: the identification / setup move. \blank{4.0cm}
  \item TODO: the computation or interpretation.\\[3pt]
        \writelines{2}
  \item \textbf{TODO: the part that is the point of the box} --- the one that tests the
        misconception on new ground.\\[3pt]
        \writelines{2}
\end{enumerate}
\end{practicebox}


% ── THE NOTES END AT GUIDED PRACTICE ─────────────────────────────────────────
% There is deliberately no "you do" box here: ../homework/ is the individual practice,
% started in class in the last 8 minutes while the teacher circulates.

\end{document}
